% !TeX root = ../main.tex
\section{Discussion}
\label{sec:discussion}

\paragraph{What a practitioner should take from this.}
If you are training a small model to do MI against an LLM oracle, three things follow.
(i)~\textbf{Do not stop at the reward curve.} Both arms here would report a large, highly
significant, multi-rubric success, and one of them ends with a therapist that asks a
question every seventh turn. (ii)~\textbf{Buy a second grader before you buy more
iterations.} Re-scoring the whole grid with a held-out judge cost less than a single
training iteration and is what separated a real gain from a grader-specific one. (iii)~
\textbf{Prefer the method that explores}, and treat iteration count as a hyperparameter with
a validation criterion attached --- GRPO's best model is four iterations before its last.

\paragraph{Why the winner wins.}
Our reading is that PTO's advantage is not the DPO loss (\S\ref{sec:mechanism} rules that
out fairly directly) but the sequential trunk. Because each branch point commits to its
winner and continues, PTO collects candidates in states reached by a run of good choices ---
states that GRPO's fixed, once-per-iteration prompt list never contains. That is a
distribution-of-training-contexts advantage, and it also predicts the observed cost: PTO's
sequential growth cannot be batched as freely, and its $\tau$ filter throws data away.

\paragraph{Why the loser loses in a specific way.}
GRPO's failure is not instability in the usual RL sense --- it improves monotonically for
eight iterations and its variance is unremarkable. It fails by \emph{succeeding at the
stated objective}: it drives the reward up, and the reward is a 22-item questionnaire in
which self-disclosure and taking charge are credited. The clinical failure mode
(a warm, praising, non-inquiring counsellor) and the machine-learning failure mode (a policy
overfitting its grader) turn out to be the same event viewed from two instruments, which is
what makes MI a good testbed for reward hacking rather than merely a domain in which it
occurs.

\paragraph{On the instruments.}
The most uncomfortable finding for the paradigm is \S\ref{sec:results}'s factor structure:
five rubrics that sound like five constructs behave like one. A study that reports Q1, Q2,
WAI-SR, CSQ-8 and MI-SAT all improving has, to a first approximation, reported one number
five times. Batteries for LLM-graded MI should be chosen for what they \emph{dis}agree
about --- which, empirically here, means MI-inconsistency counts and the MITI technique
ratios --- and not for how many validated instruments they contain.

\section{Conclusion}
\label{sec:conclusion}

We compared preference-tree and group-relative optimization for training a 1B therapist
model to do Motivational Interviewing, holding the base policy, the simulated patients, the
grading oracle, the candidate budget and the look-ahead depth fixed. Both raise the reward
substantially; PTO leads at the matched ten-iteration endpoint and keeps improving where
GRPO peaks and regresses. But the gains in both arms come with a drift toward affirmation
and away from inquiry, and a judge that never touched training credits $0.80$ of PTO's Q1
gain against $0.28$ of GRPO's --- an onset curve that begins at iteration~3. Reading the
training signal directly shows the two losses extract nearly the same direction from the
same candidates, so the method gap is about which conversational states get explored, not
about DPO versus group-relative weighting. We recommend gain retention under a decoupled
grader as a routine diagnostic wherever an LLM fills in a questionnaire for reward.
